\section{An explicit minimum-cut distance obstruction}
\label{sec:four-high-connectivity-long-cut}

The proposed short-cycle rule for minimum cyclic cuts is false, already at cut order five.
The following construction uses an explicit cubic Cayley graph.
Its only external input is a published cyclic-connectivity theorem.

\begin{proposition}
\label{prop:four-honeycomb-long-cut}
For every integer $M\ge10$, there is a connected simple cubic graph $G_M$ of rank
\[
 r_M=40M+6
\]
with a unique minimum cyclic cut of size five.
Every cycle containing an edge of that cut has length at least
\[
 2M+2=\frac{r_M-6}{20}+2.
\]
The graph has girth six.
\end{proposition}

\begin{proof}
Put $L=5M$ and $A=\mathbb Z_L\times\mathbb Z_4$.
Let $H_M$ have vertices $A\times\{0,1\}$.
For each $z\in A$, join $(z,0)$ to
\[
 (z,1),\qquad (z+(1,0),1),\qquad (z+(0,1),1).
\]
This graph is simple, cubic, and connected.
Translations of $A$ preserve its edges.
The map $(z,s)\mapsto(-z,1-s)$ also preserves its edges.
These automorphisms act transitively on the vertices.

The graph is bipartite.
The six nonzero ordered differences of $0,(1,0),(0,1)$ are distinct in $A$.
Thus two vertices in the same bipartition class have at most one common neighbor, excluding four-cycles.
The following sequence is a six-cycle:
\[
 ((0,0),0),\ ((1,0),1),\ ((1,-1),0),\
 ((1,-1),1),\ ((0,-1),0),\ ((0,0),1).
\]
Hence $H_M$ has girth six.
Theorem~17 of Nedela and \v{S}koviera identifies the cyclic connectivity of a cubic vertex-transitive graph with its girth.
Therefore $H_M$ has cyclic connectivity six.
The primary source is
\href{https://dml.cz/bitstream/handle/10338.dmlcz/132921/MathSlov_45-1995-5_3.pdf}{\emph{Atoms of cyclic connectivity in cubic graphs}, Mathematica Slovaca 45 (1995), 481--499}, page~498.

Subdivide the five edges
\[
 ((jM,0),0)((jM,0),1),\qquad 0\le j\le4,
\]
once each, with new vertices $p_0,\ldots,p_4$.
Call the resulting normalized five-pole $X_M$.
Subdivision preserves its rank and internal cyclic connectivity.
In particular,
\[
 |V(X_M)|=40M+5,\qquad \beta(X_M)=20M+1.
\]
Assign each $p_j$ the first coordinate $jM$.
Along every edge, that coordinate moves by at most one in the cycle $\mathbb Z_L$.
Distinct ports therefore have distance at least $M\ge10$.

Take two disjoint copies of $X_M$ and join corresponding ports by five edges.
Lemma~\ref{lem:random-distant-port-gluing} applies with $k=5$.
It shows that the five new edges form the unique cyclic cut of size at most five.
The order of the resulting cubic graph is $80M+10$, giving the stated rank.

Every cycle meeting the joining cut uses at least two joining edges.
It includes at least two internal segments between distinct ports.
Each segment has length at least $M$, proving the claimed lower bound.
All other cycles lie within one subdivided core and have length at least six.
An unchanged six-cycle exists away from the five subdivided edges, so the girth remains six.
\end{proof}

For sufficiently large $M$, every edge of the unique minimum cyclic cut avoids every cycle of length at most $r_M/(\log r_M)^2$.
Thus neither minimum-cut selection nor minimum-cut atoms alone can supply the proposed short cycle.
This explicit example has bounded girth.
Corollary~\ref{cor:random-long-cut-diagonal} gives a stronger qualitative obstruction using the arbitrary-fixed-cut construction.
It chooses the graph order sufficiently large at each cut order.
Both cyclic connectivity and girth then tend to infinity, while the same short-cycle rule still fails.
This does not exclude a statement with a specified growth rate for either parameter.

The exact finite checker is \texttt{four-high-connectivity-long-cut-check.py}.
It checks the graph definition, automorphisms, girth, port distances, and shortest cycles through all five joining edges.
It also checks every possible small-cut pattern after the published cyclic-connectivity reduction.
Its output is \texttt{four-high-connectivity-long-cut-checks.json}.
